\scanpage{94}
de $R$. \deleted{\unsure{Par le foncteur contravariant} \unsure{can.}}\deleted{\inserted{\unsure{identique}}} Or un ss-objet dans $R$ est un \deleted{objet \unsure{quoti}} sous-objet dans (Sch-aff), en effet $\mathbb Z[t,t^{-1}]=\underline{\mathbb Z[x,y]}/(xy-1)$. Donc O.K. \deleted{\unsure{si }\unsure{$\exists$} \unsure{lim}} avec $\varprojlim$ finies

Ceci dit, considérons pour un objet $A$ le morphisme $\deleted{\unsure{e\amalg{}}}A^\ast\longrightarrow A$, où $e$ est l'objet final. On pourra dire que $A$ (\deleted{\unsure{dans}}\inserted{anneau} dans une catégorie \deleted{quelconque}\inserted{avec $\varprojlim$ finies et $\varinjlim$ finies}) est un objet corps si ce morphisme est un isom.

Plaçons nous p.ex. dans la théorie des topos. Alors \deleted{\unsure{considér}} le topos universel \inserted{(classifiant)} pour les anneaux est $\deleted{\unsure{S}}\widehat R$ ($R\simeq\mathrm{Sch\,aff\ tf}$). L'objet anneau universel est $O=\mathbb E^1_{\mathbb Z}=\unsure{\mathrm{Spec}|\mathrm{Sp}}(\mathbb Z[t])$.

Il faut donc rendre inversible la flèche
\[
\deleted{\unsure{\mathbb Z[t]\longrightarrow B}}\qquad
\deleted{\unsure{e}}\amalg O^\ast\mathbin{\unsure{-|\to}}O
\]
(\deleted{\unsure{i.e.}}\inserted{ce n'est pas} la \unsure{transposée} (dans $R$, $D\to O$) de la flèche \inserted{(dans $\widehat R$, \uline{pas} dans $R$ !)} déjà utilisée plus haut\inserted{, car ce n'est pas une flèche dans $R$ !}) On voit que cela implique que \deleted{tte flèche $X\mathbin{-}Y$ \inserted{dans} $R$ qui est un morphisme bijectif de schémas : ext. résiduelles triviales devient un isom dans le nouveau topos classifiant $B'$. [ceci encore en remarquant que} pour tt \inserted{$X\in\widehat R$ et tt} $f\in\Gamma(X,O_X)$,
\[
 X_f\amalg V(f)\mathbin{-}X
\]
doit devenir un isom, car déduit de la situation précédente par chgt de base
\diagram{\begin{tikzcd}[ampersand replacement=\&,row sep=small]
 X\arrow[r,"f"] \& O=\mathbb E^1_{\mathbb Z}\\
 \& O^\ast\amalg e\arrow[u]
\end{tikzcd}}
] \deleted{Or la catégorie quotient de $R$ par l'inv des flèches précédentes (qui manifestement est stable par $\varprojlim$) est la catégorie des schémas \unsure{compacts|composés|complet}}
% archive sheet number: 93
